\documentclass[11pt,twocolumn]{article}
\usepackage[margin=1in]{geometry}
\usepackage{amsmath,amssymb,amsthm}
\usepackage{hyperref}
\usepackage{cite}
\usepackage{graphicx}
\usepackage{algorithm}
\usepackage{algpseudocode}
\usepackage{booktabs}
\usepackage{microtype}

\newtheorem{theorem}{Theorem}
\newtheorem{lemma}{Lemma}
\newtheorem{corollary}{Corollary}
\newtheorem{definition}{Definition}

\title{\textbf{nunc}: Proving Wall-Clock Time via\\
Multi-Source Network Consensus with\\
Cryptographic and Statistical Guarantees}

\author{
  Nick Spiker \\
  \texttt{holdmyoscilloscope.com}
}

\date{}

\begin{document}
\raggedright

\maketitle

\begin{abstract}
I present \textsc{nunc}, a protocol for establishing trustworthy wall-clock
time using the ambient infrastructure of the public internet — without
dedicated time-stamping servers, without spending cryptocurrency, and without
trusting any single party.  The protocol selects a TRNG-seeded random subset
of a large, diverse pool of independently-operated servers, verifies their
responses against hardcoded public keys anchored to Certificate Transparency
(CT) logs, and computes a consensus timestamp via interval intersection on the
agreeing majority.  I reduce the correctness argument to Lamport and
Melliar-Smith's Byzantine clock synchronization result~\cite{lamport1984}, show
that TRNG unpredictability combined with pool diversity makes Byzantine
compromise of a selection majority computationally infeasible, and demonstrate
that the organic timestamp distribution admits a Kolmogorov-Smirnov
goodness-of-fit test that provides a continuous, real-valued integrity signal —
detecting coordinated manipulation as a statistically anomalous departure from
the expected unimodal blackbody-like curve.  The result is the first protocol
to unify Byzantine fault tolerance, cryptographic certificate provenance, and
statistical distribution analysis into a single time-establishment proof over
open, heterogeneous internet infrastructure.
\end{abstract}

%------------------------------------------------------------------
\section{Introduction}
%------------------------------------------------------------------

Every security property that depends on time --- token expiry, certificate
validity windows, replay-attack detection --- rests on an assumption that is
almost universally unverified: that the local system clock is correct.

The standard remedy is the Network Time Protocol (NTP)~\cite{rfc5905}, which
is unauthenticated by default and provides no mechanism for a client to detect
a lying server.  Network Time Security (NTS)~\cite{rfc8915} adds
authentication but still requires trust in a small set of designated servers.
Roughtime~\cite{roughtime} introduces nonce-binding and malfeasance proofs but,
as of this writing, has fewer than five publicly reachable servers, rendering
its ecosystem-level guarantees nominal.

I observe that the public internet already contains millions of independently
operated servers that broadcast their local time in standard protocol headers:
HTTPS \texttt{Date:} fields, NTP packets, SMTP banners, and others.  Each
server is independently operated, independently clocked, and in aggregate spans
every jurisdiction, autonomous system, and legal domain on earth.  No
coordination among them is required.  No new infrastructure is needed.

\textsc{nunc} harvests this ambient signal.  The core insight is that
combining (1) unpredictable server selection via a true random nonce, (2)
cryptographic binding of responses to hardcoded public keys via CT logs, and
(3) a statistical goodness-of-fit test on the resulting timestamp distribution
yields a proof of time whose assumptions reduce entirely to standard
cryptographic primitives and 40-year-old distributed systems theory.

%------------------------------------------------------------------
\section{Background and Related Work}
%------------------------------------------------------------------

\subsection{Byzantine Clock Synchronization}

Lamport and Melliar-Smith~\cite{lamport1984} proved that a set of $n$ clocks
can be synchronized in the presence of up to $f$ Byzantine (arbitrarily
faulty) clocks if and only if $n > 3f$.  Their Interactive Convergence
Algorithm takes the median of received clock values, discarding outliers, and
provably converges to a synchronized state within a bounded error that depends
only on message delay uncertainty and clock drift rate.  This result was
subsequently machine-checked at SRI~\cite{rushby1991} and extended to the
self-stabilizing setting~\cite{daliot2003}.

\textsc{nunc} does not re-prove this result.  I reduce to it.

\subsection{Certificate Transparency}

Certificate Transparency~\cite{rfc6962} requires that every publicly trusted
TLS certificate be submitted to at least one append-only, publicly auditable
log before browsers will accept it.  Each log entry includes a Signed
Certificate Timestamp (SCT) — a signature by the log's private key over the
certificate hash and log timestamp.  The logs are Merkle hash chains; retroactive
insertion of a fraudulent entry requires rewriting the chain from that point,
which is immediately detectable by any party that has observed the prior log
head.

I use CT log public keys as compile-time trust anchors, eliminating the need
to trust any runtime certificate authority.

\subsection{Roughtime}

Roughtime~\cite{roughtime} binds each server response to a client-supplied
nonce, making replay attacks detectable, and chains responses across servers so
that any lying server produces a signed, irrefutable record of its
malfeasance.  The protocol inherits the nonce-binding intuition but does not
require Roughtime infrastructure; I achieve analogous freshness guarantees via
TRNG-seeded selection and SCT verification.

%------------------------------------------------------------------
\section{Protocol}
%------------------------------------------------------------------

\subsection{Definitions}

\begin{definition}[Pool]
  A pool $P = \{(h_i, pk_i, \mathit{asn}_i, \mathit{cc}_i)\}_{i=1}^{N}$ is a
  set of server entries, each comprising a hostname/IP $h_i$, a hardcoded
  public key $pk_i$, an autonomous system number $\mathit{asn}_i$, and a
  country code $\mathit{cc}_i$.  Public keys are pinned at compile time and
  anchored to CT log entries.
\end{definition}

\begin{definition}[Nonce]
  A nonce $\eta \leftarrow \{0,1\}^{64}$ is drawn from a hardware true random
  number generator (TRNG) immediately before each query.
\end{definition}

\begin{definition}[Selection]
  $S(\eta, P, k)$ selects $k$ entries from $P$ deterministically given $\eta$,
  maximizing ASN and country-code diversity using a seeded shuffle followed by
  a greedy diversity pass.
\end{definition}

\begin{definition}[Observation]
  For each selected server $i$, the client records an observation
  $o_i = (t_i, \rho_i, \sigma_i)$ where $t_i$ is the reported timestamp,
  $\rho_i$ is the measured round-trip time, and $\sigma_i$ is the TLS SCT
  signature.
\end{definition}

\subsection{Algorithm}

\begin{algorithm}
\caption{\textsc{nunc} time establishment}
\begin{algorithmic}[1]
  \State $\eta \leftarrow \mathrm{TRNG}()$
  \State $Q \leftarrow S(\eta, P, k)$
  \State Issue parallel HEAD requests to all $q \in Q$
  \For{each response $(t_i, \rho_i, \sigma_i)$}
    \If{$\neg\,\mathrm{VerifySCT}(\sigma_i, pk_i)$} discard $o_i$ \EndIf
  \EndFor
  \State $\hat{t} \leftarrow \mathrm{median}(\{t_i\})$
  \State Reject $o_i$ if $|t_i - \hat{t}| > \Delta_{\mathrm{reject}}$
  \State $[L, R] \leftarrow \bigcap_i [t_i - \rho_i/2,\; t_i + \rho_i/2]$
  \State $T^* \leftarrow (L+R)/2$, $\;\epsilon \leftarrow (R-L)/2$
  \State $p \leftarrow \mathrm{KS\text{-}test}(\{t_i\}, \hat{F})$
  \State \Return $(T^*, \epsilon, p)$
\end{algorithmic}
\end{algorithm}

The client returns the consensus midpoint $T^*$, a confidence half-width
$\epsilon$, and a KS p-value $p$ indicating distributional integrity.

%------------------------------------------------------------------
\section{Correctness}
%------------------------------------------------------------------

\subsection{Reduction to Lamport}

\begin{theorem}[Correctness under honest majority]
  If fewer than $k/3$ of the $k$ selected servers are Byzantine, then $T^*$
  is within the Lamport error bound of true time $T$.
\end{theorem}

\begin{proof}
  Honest server $i$ reports $t_i$ satisfying $|t_i - T| \leq \varepsilon_i$
  where $\varepsilon_i$ is bounded by its upstream clock accuracy (GPS or
  atomic stratum 1, typically $<1$ms).  The interval intersection of honest
  observations contains $T$.  By Lamport and
  Melliar-Smith~\cite{lamport1984}, the median-based convergence algorithm
  with fewer than $k/3$ faulty inputs produces a value within the proven
  error bound.  The outlier rejection step is precisely the Interactive
  Convergence Algorithm's discard rule.  $\square$
\end{proof}

\subsection{Honest Majority from TRNG and Pool Diversity}

\begin{lemma}[Selection unpredictability]
  An adversary who does not know $\eta$ prior to its generation cannot
  predict $S(\eta, P, k)$ with probability greater than $k/N$.
\end{lemma}

\begin{proof}
  $\eta$ is drawn uniformly from $\{0,1\}^{64}$ by a hardware TRNG.
  $S(\eta, P, k)$ is a deterministic function of $\eta$.  Without $\eta$,
  the adversary's best strategy is to guess; the probability of any specific
  server being selected is $k/N$.  $\square$
\end{proof}

\begin{lemma}[Pre-compromise infeasibility]
  The probability that an adversary controls $\geq k/3$ of $S(\eta, P, k)$
  by pre-compromising $m$ servers in $P$ is at most
  $\binom{m}{k/3}\binom{N-m}{2k/3} / \binom{N}{k}$.
\end{lemma}

\begin{proof}
  This is the hypergeometric tail probability.  For the reference implementation
  pool of $N = 6{,}732$ entries with $k = 42$ and $m = 67$ (adversary controls
  1\% of the pool), the probability of $\geq 14$ compromised servers appearing
  in the selection is $\approx 9.5 \times 10^{-19} < 10^{-18}$.  The bound
  scales favourably with $N$: for fixed adversary fraction $m/N$, increasing
  $N$ strictly decreases this probability.  $\square$
\end{proof}

\begin{theorem}[Honest majority with overwhelming probability]
  Under TRNG unpredictability and pool diversity, the selected set has an
  honest majority except with negligible probability in the security parameter.
\end{theorem}

\begin{proof}
  Combine Lemmas 1 and 2.  The adversary cannot predict the selection
  (Lemma 1), and pre-compromising a majority of the pool large enough to
  guarantee $\geq k/3$ selected faulty servers requires compromising $O(N/3)$
  independently-operated servers across independent ASNs, jurisdictions, and
  CT log domains — a task whose cost grows linearly with pool size and is
  computationally infeasible for any realistic $N$.  $\square$
\end{proof}

\subsection{Cryptographic Binding via CT Logs}

\begin{theorem}[Response authenticity]
  A response bearing a valid SCT against a hardcoded CT log public key
  originated from the genuine private key holder, except under collision
  resistance of the underlying hash function.
\end{theorem}

\begin{proof}
  An SCT is a signature by the CT log's private key over the certificate's
  hash and issuance timestamp.  The certificate itself binds the server's
  public key.  Forging a valid SCT requires either breaking the log's
  signature scheme, finding a hash collision in the Merkle tree, or
  compromising the log's private key.  Each CT log is independently operated;
  the adversary must compromise all logs that have signed the certificate.
  Under standard cryptographic assumptions (EUF-CMA security of ECDSA/Ed25519,
  collision resistance of SHA-256), this is computationally infeasible.
  $\square$
\end{proof}

%------------------------------------------------------------------
\section{Statistical Integrity Signal}
%------------------------------------------------------------------

\subsection{The Expected Distribution}

Honest server timestamps are independently drawn from clocks that track true
time to within $\varepsilon_i \ll 1$s.  Network jitter adds symmetric noise
$\sim \mathcal{N}(0, \sigma^2_{\rho})$.  CDN staleness adds a right-skewed
component with exponential decay.  The resulting distribution is unimodal with
a tight peak near $T$ and a long right tail — qualitatively similar in shape
to a blackbody spectral curve.

Let $\hat{F}$ denote the empirical CDF of this distribution, established from
a baseline of clean measurements.

\subsection{Kolmogorov-Smirnov Integrity Test}

\begin{definition}[KS statistic]
  Given observations $\{t_i\}_{i=1}^k$ and reference CDF $\hat{F}$, the KS
  statistic is:
  \[
    D_k = \sup_t \left| F_k(t) - \hat{F}(t) \right|
  \]
  where $F_k$ is the empirical CDF of $\{t_i\}$.
\end{definition}

\begin{theorem}[Manipulation detectability]
  A coordinated attack injecting $m$ false timestamps clustered at $T' \neq T$
  produces a bimodal distribution whose KS statistic $D_k$ exceeds the
  $(1-\alpha)$ critical value with probability approaching 1 as
  $m/k \to 1/3$.
\end{theorem}

\begin{proof}
  A bimodal distribution with modes at $T$ and $T'$ has an empirical CDF that
  departs from the unimodal $\hat{F}$ by approximately $(m/k)$ at the
  midpoint $(T+T')/2$.  By the Glivenko-Cantelli theorem, $D_k \to m/k$
  almost surely as $k \to \infty$.  For $m/k \geq 1/3$, this departure
  exceeds the KS critical value at any standard significance level, yielding
  detection with overwhelming probability.  $\square$
\end{proof}

\subsection{Continuous Integrity Monitoring}

The KS p-value $p$ returned by \textsc{nunc} is a continuous real-valued
integrity signal:

\begin{itemize}
  \item $p \approx 1$: distribution consistent with honest baseline.
  \item $p \to 0$: anomalous distribution — active manipulation, network
        degradation, or infrastructure incident.
  \item Monotone drift in $p$ over successive measurements: slow probing or
        gradual infrastructure compromise.
\end{itemize}

This converts time establishment from a binary trusted/untrusted judgment into
a continuously monitored integrity metric with a formal statistical backing.

%------------------------------------------------------------------
\section{Main Result}
%------------------------------------------------------------------

\begin{theorem}[\textsc{nunc} correctness]
  Under the assumptions of (i) TRNG unpredictability, (ii) EUF-CMA security
  of SCT signature schemes, (iii) collision resistance of CT log hash
  functions, and (iv) honest upstream clocks at stratum 1 sources, the
  \textsc{nunc} protocol returns $T^*$ satisfying
  \[
    |T^* - T| \leq \frac{R - L}{2} + \varepsilon_{\max}
  \]
  with probability $1 - \mathit{negl}(\lambda)$, where $[L, R]$ is the
  consensus interval, $\varepsilon_{\max}$ is the maximum honest server
  clock error, and $\lambda$ is the security parameter.  Furthermore, any
  coordinated manipulation of $\geq k/3$ sources is detected by the KS test
  with probability approaching 1.
\end{theorem}

\begin{proof}
  By Theorem 3, the selected set has an honest majority except with negligible
  probability.  By Theorem 4, honest responses are cryptographically
  authenticated.  By Theorem 1 (reduction to Lamport), an honest majority
  yields $T^*$ within the stated bound.  By Theorem 5, minority manipulation
  is statistically detectable.  $\square$
\end{proof}

%------------------------------------------------------------------
\section{Discussion}
%------------------------------------------------------------------

\subsection{Threat Model Ceiling}

The construction does not defend against an adversary who controls the
client's physical uplink at the BGP level.  Such an adversary can drop all
responses, forcing a timeout.  However, they cannot forge a valid consensus
— authenticated silence is detectable, authenticated lies are not feasible.
The protocol fails safely: no result rather than a wrong result.

\subsection{Pool Maintenance}

The security argument scales with pool size $N$.  The reference implementation
ships with $N = 6{,}732$ vetted entries spanning 206 countries and 6 protocols;
this is sufficient for the hypergeometric bound to yield $< 10^{-18}$ at 1\%
adversary control.  Pool entries must be periodically validated (live
connection test, SCT verification, ASN check) and the server list refreshed.
Attrition in a diverse pool of this size is operationally manageable, and
growing $N$ monotonically strengthens the security guarantee.

\subsection{Relationship to Roughtime}

Roughtime provides a stronger per-response freshness guarantee via explicit
nonce binding.  \textsc{nunc} achieves analogous freshness through TRNG
selection unpredictability.  As the Roughtime ecosystem matures, its servers
are natural high-weight additions to the \textsc{nunc} pool.

%------------------------------------------------------------------
\section{Conclusion}
%------------------------------------------------------------------

I have presented \textsc{nunc}, a protocol for proving wall-clock time over
ambient internet infrastructure.  The proof reduces to Lamport's 40-year-old
Byzantine clock synchronization result, stands on standard cryptographic
assumptions via CT log anchoring, and introduces a novel continuous integrity
signal via Kolmogorov-Smirnov goodness-of-fit testing against the expected
unimodal timestamp distribution.  The protocol requires no dedicated
infrastructure, no cryptocurrency, and no trust in any single party.  An open
source reference implementation is available as the \texttt{nunc-time} Rust
crate.

\begin{thebibliography}{9}

\bibitem{lamport1984}
  L.~Lamport and P.~M.~Melliar-Smith,
  ``Byzantine clock synchronization,''
  in \textit{Proc. 3rd ACM Symp. Principles of Distributed Computing},
  pp.~68--74, 1984.

\bibitem{rushby1991}
  J.~Rushby and F.~von Henke,
  ``Formal verification of the interactive convergence clock synchronization
  algorithm,''
  SRI-CSL-89-3R, SRI International, 1989 (revised 1991).

\bibitem{daliot2003}
  A.~Daliot, D.~Dolev, and H.~Parnas,
  ``Self-stabilizing pulse synchronization inspired by biological pacemaker
  networks,''
  in \textit{Proc. 6th Symp. Self-Stabilizing Systems}, 2003.

\bibitem{rfc5905}
  D.~Mills et al.,
  ``Network Time Protocol Version 4: Protocol and Algorithms Specification,''
  RFC~5905, IETF, 2010.

\bibitem{rfc8915}
  D.~Franke et al.,
  ``Network Time Security for the Network Time Protocol,''
  RFC~8915, IETF, 2020.

\bibitem{rfc6962}
  B.~Laurie, A.~Langley, and E.~Kasper,
  ``Certificate Transparency,''
  RFC~6962, IETF, 2013.

\bibitem{roughtime}
  A.~Langley, A.~Malhotra, and W.~Ladd,
  ``Roughtime,''
  IETF Internet-Draft \textit{draft-ietf-ntp-roughtime-19}, 2026.

\end{thebibliography}

\end{document}